\documentclass[12pt, a4paper]{article}
\usepackage[spanish]{babel}
\usepackage[utf8]{inputenc}
\usepackage{amsmath, amssymb, amsthm}
\usepackage{mathtools}
\usepackage{stmaryrd}   % \llbracket, \rrbracket
\usepackage{hyperref}
\usepackage{geometry}
\geometry{margin=2.5cm}

\newtheorem{definition}{Definición}[section]
\newtheorem{theorem}{Teorema}[section]
\newtheorem{lemma}{Lema}[section]
\newtheorem{proposition}{Proposición}[section]
\newtheorem{remark}{Observación}[section]
\newtheorem{example}{Ejemplo}[section]

\title{Formalización de la Destrucción de SSA en CVM\\[0.5em]
\large Borrador inicial --- Definiciones previas}
\author{}
\date{}

\begin{document}
\maketitle
\tableofcontents
\newpage

% =========================================================
\section{Introducción}
% =========================================================

La forma estática de asignación única (\emph{Static Single Assignment}, SSA) es una
representación intermedia del código en la que cada variable es definida exactamente una vez.
Esta propiedad simplifica numerosos análisis y optimizaciones de compiladores. Sin embargo, SSA
no es directamente ejecutable: las \emph{funciones phi} ($\phi$-funciones) que aparecen en los
puntos de convergencia del flujo de control no tienen correspondencia directa en bytecode
convencional.

El proceso de \emph{destrucción de SSA} consiste en eliminar las $\phi$-funciones y obtener un
programa equivalente sin ellas. En este trabajo formalizamos la destrucción \emph{naive}
implementada para el bytecode CVM: una versión sin optimizar que, pese a poder introducir copias
redundantes, garantiza la corrección mediante el uso de variables temporales frescas.

El objetivo de este documento es establecer las definiciones matemáticas necesarias para
demostrar formalmente que el código CVM antes y después de pasar por SSA (construir y destruir)
produce el mismo comportamiento observable. La corrección de la \emph{construcción} de SSA se
asume como hipótesis; aquí únicamente nos interesa la \emph{destrucción}.

% =========================================================
\section{El lenguaje CVM-CFG}
% =========================================================

\subsection{Variables y valores}

\begin{definition}[Variables y valores]
Sea $\mathcal{V}$ un conjunto numerable de \emph{nombres de variable}. Sea $\mathcal{D}$ el
\emph{dominio de valores}, que comprende enteros de 64 bits ($\mathbb{Z}_{2^{64}}$) y elementos
de cuerpos finitos. Una \emph{valuación} es una función parcial
$\sigma : \mathcal{V} \rightharpoonup \mathcal{D}$.
\end{definition}

\begin{definition}[Expresiones atómicas]
Una \emph{expresión atómica} $a$ es uno de los siguientes:
\begin{itemize}
  \item Una constante $c \in \mathcal{D}$,
  \item Una variable $x \in \mathcal{V}$.
\end{itemize}
El conjunto de expresiones atómicas se denota $\mathcal{A}$.
\end{definition}

\begin{definition}[Evaluación de una expresión atómica]
Dada una valuación $\sigma$, la evaluación de $a \in \mathcal{A}$ se define como:
\[
  \llbracket a \rrbracket_\sigma =
  \begin{cases}
    c         & \text{si } a = c \in \mathcal{D},\\
    \sigma(x) & \text{si } a = x \in \mathcal{V}.
  \end{cases}
\]
\end{definition}

\begin{definition}[Operadores]
Sea $\Omega$ un conjunto finito de \emph{operadores} (suma, resta, multiplicación, comparaciones,
operaciones de bits, etc.). Cada operador $\omega \in \Omega$ tiene una \emph{aridad}
$\mathrm{ar}(\omega) \geq 0$ y una \emph{semántica}
$\llbracket \omega \rrbracket : \mathcal{D}^{\mathrm{ar}(\omega)} \to \mathcal{D}$.
\end{definition}

\subsection{Instrucciones}

\begin{definition}[Instrucción]
\label{def:instruccion}
Una \emph{instrucción} es una tupla $s = (\tau,\, x,\, \omega,\, \vec{a})$ donde:
\begin{itemize}
  \item $\tau \in \{\texttt{i64},\, \texttt{ff},\, \bot\}$ es el tipo numérico (opcional),
  \item $x \in \mathcal{V} \cup \{\bot\}$ es la variable de salida (opcional),
  \item $\omega \in \Omega \cup \{\bot\}$ es el operador (opcional),
  \item $\vec{a} = (a_1, \ldots, a_n)$ es el vector de operandos, con $a_i \in \mathcal{A}$.
\end{itemize}
Cuando $x \neq \bot$, $\omega = \bot$ y $n = 1$, la instrucción se denomina
\emph{instrucción de copia} y se escribe $x \leftarrow a_1$.
\end{definition}

\begin{definition}[Ejecución de una instrucción]
\label{def:ejec-instruccion}
La ejecución de $s = (\tau, x, \omega, \vec{a})$ sobre la valuación $\sigma$ produce:
\[
  \llbracket s \rrbracket_\sigma =
  \begin{cases}
    \sigma\bigl[x \mapsto \llbracket \omega \rrbracket\bigl(\llbracket a_1 \rrbracket_\sigma, \ldots, \llbracket a_n \rrbracket_\sigma\bigr)\bigr]
      & \text{si } x \neq \bot \text{ y } \omega \neq \bot,\\[4pt]
    \sigma\bigl[x \mapsto \llbracket a_1 \rrbracket_\sigma\bigr]
      & \text{si } x \neq \bot \text{ y } \omega = \bot \text{ (copia)},\\[4pt]
    \sigma
      & \text{si } x = \bot \text{ (sin variable de salida)}.
  \end{cases}
\]
Aquí $\sigma[x \mapsto v]$ denota la valuación que coincide con $\sigma$ en toda variable
excepto $x$, a la que asigna el valor $v$.
\end{definition}

\begin{definition}[Ejecución de una secuencia de instrucciones]
\label{def:ejec-secuencia}
La ejecución de la secuencia $S = [s_1, \ldots, s_k]$ sobre $\sigma$ se define inductivamente:
\[
  \llbracket [\,] \rrbracket_\sigma = \sigma, \qquad
  \llbracket [s_1, \ldots, s_k] \rrbracket_\sigma =
    \llbracket [s_2, \ldots, s_k] \rrbracket_{\llbracket s_1 \rrbracket_\sigma}.
\]
Las instrucciones se ejecutan \textbf{secuencialmente}: cada instrucción ve el estado modificado
por todas las anteriores.
\end{definition}

\subsection{Bloques básicos y grafo de flujo de control}

\begin{definition}[Sucesor de un bloque]
El \emph{sucesor} de un bloque básico es uno de los siguientes:
\begin{itemize}
  \item \textbf{Incondicional}: $\mathrm{uncond}(j)$, que transfiere el control al bloque $j$.
  \item \textbf{Condicional}: $\mathrm{cond}(a, j_1, j_2)$, con $a \in \mathcal{A}$. Si
    $\llbracket a \rrbracket_\sigma \neq 0$, el control pasa a $j_1$; en caso contrario, a $j_2$.
  \item \textbf{Terminal}: $\bot$ (bloque final, sin sucesor).
\end{itemize}
\end{definition}

\begin{definition}[Bloque básico]
\label{def:bloque}
Un \emph{bloque básico} es una tupla
$B = (\mathrm{id},\, \Phi,\, S,\, \mathrm{pred},\, \mathrm{succ})$ donde:
\begin{itemize}
  \item $\mathrm{id} \in \mathbb{N}$ es el identificador único del bloque,
  \item $\Phi = [\phi_1, \ldots, \phi_m]$ es la lista (posiblemente vacía) de $\phi$-funciones,
  \item $S = [s_1, \ldots, s_k]$ es la secuencia de instrucciones,
  \item $\mathrm{pred} \subseteq \mathbb{N}$ es el conjunto de identificadores de predecesores,
  \item $\mathrm{succ}$ es el sucesor del bloque.
\end{itemize}
\end{definition}

\begin{definition}[Grafo de flujo de control (CFG)]
\label{def:cfg}
Un \emph{grafo de flujo de control} (CFG) es un par $G = (\mathcal{B},\, \mathrm{entry})$ donde:
\begin{itemize}
  \item $\mathcal{B} = [B_0, B_1, \ldots, B_{n-1}]$ es una lista indexada de bloques básicos,
  \item $\mathrm{entry} \in \mathbb{N}$ es el identificador del bloque de entrada.
\end{itemize}
La relación de aristas del grafo está determinada por los sucesores: existe una arista $(i, j)$
si y solo si $j$ aparece en $\mathrm{succ}(B_i)$.

Dado un bloque $B_j$, sus \emph{predecesores} son $\mathrm{pred}(B_j) = \{i \mid (i,j) \in E\}$,
consistente con el campo $\mathrm{pred}$ de la Definición~\ref{def:bloque}.
\end{definition}

\begin{definition}[Programa CVM (forma no-SSA)]
Un \emph{programa CVM} es un CFG $G = (\mathcal{B},\, \mathrm{entry})$ en el que todos los
bloques tienen $\Phi = [\,]$ (sin $\phi$-funciones). Las instrucciones pueden asignar múltiples
veces a la misma variable.
\end{definition}

% =========================================================
\section{Forma SSA}
% =========================================================

\subsection{La propiedad SSA}

\begin{definition}[Forma SSA]
\label{def:ssa}
Un CFG $G = (\mathcal{B},\, \mathrm{entry})$ está en \emph{forma SSA} si cada variable
$x \in \mathcal{V}$ aparece como variable de salida en \emph{a lo sumo una} instrucción o
$\phi$-función en todo el grafo.

Formalmente, para todo $x \in \mathcal{V}$:
\[
  \bigl|\{B_j \in \mathcal{B} \mid \exists\, s \in S_j : \mathrm{out}(s) = x\}\bigr|
  +
  \bigl|\{B_j \in \mathcal{B} \mid \exists\, \phi \in \Phi_j : \mathrm{out}(\phi) = x\}\bigr|
  \leq 1,
\]
donde $\mathrm{out}(s)$ (resp. $\mathrm{out}(\phi)$) denota la variable de salida de $s$ (resp.
$\phi$), si existe.
\end{definition}

\begin{remark}
En la implementación, las variables SSA se nombran $v000, v001, v002, \ldots$ y provienen de
renombrar las variables originales del programa CVM (que pueden ser $x$, $y$, $z$, etc.). Cada
nueva definición de una variable original genera un nombre SSA fresco. La función
$\rho : \mathcal{V}_{\mathrm{SSA}} \to \mathcal{V}_{\mathrm{orig}}$ que asigna a cada nombre SSA
su nombre original se construye durante la fase de renombramiento.
\end{remark}

\subsection{Funciones phi}

\begin{definition}[$\phi$-función]
\label{def:phi}
Una \emph{$\phi$-función} en un bloque $B_j$ con predecesores $\mathrm{pred}(B_j) = \{p_1,\ldots,p_r\}$ es una tupla
\[
  \phi = \Bigl(x,\; \bigl[(y_1, p_1),\, \ldots,\, (y_r, p_r)\bigr]\Bigr)
\]
donde $x \in \mathcal{V}_{\mathrm{SSA}}$ es la variable definida y $(y_i, p_i)$ indica que si el
control llega a $B_j$ desde el predecesor $p_i$, entonces $x$ toma el valor de la variable $y_i$.

Usaremos la notación $\phi.\mathrm{out} = x$ y $\phi.\mathrm{arg}(p_i) = y_i$.
\end{definition}

\begin{remark}
Las $\phi$-funciones se sitúan siempre al \emph{inicio} del bloque, antes de cualquier
instrucción ordinaria. Esta posición es convencional: en la semántica, las $\phi$-funciones de
un bloque se ejecutan \emph{simultáneamente} al entrar en él, con la semántica de copia paralela
que se define a continuación.
\end{remark}

\subsection{Semántica de copia paralela}

La semántica fundamental de las $\phi$-funciones es que todas se evalúan \textbf{al mismo tiempo}:
los valores fuente se leen en la valuación \emph{previa a cualquier modificación}, y solo
entonces se escriben los destinos. Esto es lo que se llama \emph{semántica de copia paralela},
en contraste con la \emph{semántica secuencial} de las instrucciones ordinarias.

\begin{definition}[Ejecución paralela de $\phi$-funciones]
\label{def:ejec-phi}
Sea $B_j$ un bloque con $\Phi_j = [\phi_1, \ldots, \phi_m]$, donde
$\phi_i = (x_i, [(y_i^{p_1}, p_1), \ldots, (y_i^{p_r}, p_r)])$.
Sea $p \in \mathrm{pred}(B_j)$ el bloque predecesor desde el que se llega a $B_j$.

La \emph{ejecución paralela} de $\Phi_j$ con predecesor $p$ sobre la valuación $\sigma$ produce:
\[
  \llbracket \Phi_j \rrbracket_{\sigma,\, p}
  = \sigma\Bigl[x_1 \mapsto \sigma(y_1^p),\; x_2 \mapsto \sigma(y_2^p),\; \ldots,\; x_m \mapsto \sigma(y_m^p)\Bigr]
\]
donde \textbf{todos} los valores $\sigma(y_i^p)$ se evalúan en la valuación \emph{original}
$\sigma$, antes de modificar ningún $x_i$.
\end{definition}

\begin{remark}
La diferencia con la ejecución secuencial es crucial. Compárese con la ejecución secuencial de
las asignaciones $x_1 \leftarrow y_1^p;\; x_2 \leftarrow y_2^p;\; \ldots$: en la ejecución
secuencial, $y_2^p$ se evaluaría en la valuación que ya incluye la modificación de $x_1$. En la
ejecución paralela, ambos lados se evalúan en $\sigma$.
\end{remark}

% =========================================================
\section{Semántica operacional del CFG}
% =========================================================

\subsection{Estado de ejecución y trazas}

\begin{definition}[Configuración de ejecución]
\label{def:configuracion}
Una \emph{configuración de ejecución} del CFG $G$ es un triple
$(\sigma,\, j,\, p)$ donde:
\begin{itemize}
  \item $\sigma : \mathcal{V} \rightharpoonup \mathcal{D}$ es la valuación actual,
  \item $j \in \mathbb{N}$ es el identificador del bloque que se va a ejecutar a continuación,
  \item $p \in \mathbb{N} \cup \{\bot\}$ es el identificador del bloque predecesor inmediato
    ($p = \bot$ si $j = \mathrm{entry}$).
\end{itemize}
\end{definition}

\begin{definition}[Siguiente configuración]
\label{def:siguiente}
Sea $B_j = (\mathrm{id},\, \Phi_j,\, S_j,\, \mathrm{pred}_j,\, \mathrm{succ}_j)$ un bloque y
$(\sigma, j, p)$ una configuración. La \emph{ejecución de $B_j$} transforma la configuración
del siguiente modo:
\begin{enumerate}
  \item \textbf{Ejecución paralela de $\phi$-funciones:}
    $\sigma_1 = \llbracket \Phi_j \rrbracket_{\sigma,\, p}$
    \quad (con $\sigma_1 = \sigma$ si $\Phi_j = [\,]$ o $p = \bot$).
  \item \textbf{Ejecución secuencial de instrucciones:}
    $\sigma_2 = \llbracket S_j \rrbracket_{\sigma_1}$.
  \item \textbf{Determinación del siguiente bloque:}
    \[
      \mathrm{next}(j, \sigma_2) =
      \begin{cases}
        j' & \text{si } \mathrm{succ}_j = \mathrm{uncond}(j'),\\[3pt]
        j_1 & \text{si } \mathrm{succ}_j = \mathrm{cond}(a, j_1, j_2)
              \text{ y } \llbracket a \rrbracket_{\sigma_2} \neq 0,\\[3pt]
        j_2 & \text{si } \mathrm{succ}_j = \mathrm{cond}(a, j_1, j_2)
              \text{ y } \llbracket a \rrbracket_{\sigma_2} = 0.
      \end{cases}
    \]
\end{enumerate}
La configuración resultante es $(\sigma_2,\; \mathrm{next}(j, \sigma_2),\; j)$.
\end{definition}

\begin{definition}[Traza de ejecución]
\label{def:traza}
Una \emph{traza de ejecución} del CFG $G = (\mathcal{B}, \mathrm{entry})$ a partir de la
valuación inicial $\sigma_0$ es una secuencia (finita o infinita) de configuraciones:
\[
  \pi = (\sigma_0,\, \mathrm{entry},\, \bot) \;\to\;
        (\sigma_1,\, j_1,\, \mathrm{entry}) \;\to\;
        (\sigma_2,\, j_2,\, j_1) \;\to\; \cdots
\]
donde cada transición $(\sigma_n,\, j_n,\, p_n) \to (\sigma_{n+1},\, j_{n+1},\, p_{n+1})$
se obtiene aplicando la Definición~\ref{def:siguiente} al bloque $B_{j_n}$.

La secuencia de bloques $j_0 = \mathrm{entry},\, j_1,\, j_2,\, \ldots$ se denomina el
\emph{camino de ejecución} de $\pi$.
\end{definition}

% =========================================================
\section{El algoritmo de destrucción de SSA naive}
% =========================================================

\subsection{Motivación: semántica paralela vs.\ secuencial}

Las $\phi$-funciones tienen semántica de copia paralela (Definición~\ref{def:ejec-phi}), pero
el código ejecutable solo dispone de instrucciones con semántica secuencial
(Definición~\ref{def:ejec-secuencia}). El reto de la destrucción de SSA es reemplazar las
$\phi$-funciones por instrucciones de copia ordinarias sin alterar la semántica.

\begin{example}[El problema del intercambio]
\label{ex:swap}
Considérese el siguiente fragmento de código sin SSA con dos ramas que convergen:

\begin{verbatim}
// rama p1:          // rama p2:
a = 1                a = 10
b = 2                b = 20
goto Bj              goto Bj

// bloque Bj (join):
swap(a, b)   -- queremos que a tome el valor de b y viceversa
\end{verbatim}

En SSA, las asignaciones de \texttt{a} y \texttt{b} reciben nombres únicos
y el intercambio en $B_j$ se expresa con $\phi$-funciones:

\begin{verbatim}
// rama p1:          // rama p2:
a1 = 1               a2 = 10
b1 = 2               b2 = 20
goto Bj              goto Bj

// bloque Bj:
a3 = phi(b1 : p1,  b2 : p2)   -- a3 toma el valor de b
b3 = phi(a1 : p1,  a2 : p2)   -- b3 toma el valor de a
\end{verbatim}

Nótese que un bloque con $\phi$-funciones tiene necesariamente \emph{dos o más} predecesores:
con un único predecesor no hay confluencia de flujos y no se genera ninguna $\phi$-función.

La ejecución paralela de las $\phi$-funciones es correcta: ambas leen el estado
$\sigma$ \emph{antes} de que ninguna escriba, de modo que $a_3 \leftarrow \sigma(b_1)$
y $b_3 \leftarrow \sigma(a_1)$ simultáneamente.

Si se intenta sustituir naïvemente las $\phi$-funciones por copias secuenciales
al final del predecesor $p_1$:
\[
  a_3 \leftarrow b_1;\quad b_3 \leftarrow a_1,
\]
el resultado es incorrecto: en este caso concreto las copias son independientes
(los nombres $a_3$, $b_3$, $a_1$, $b_1$ son todos distintos) y no hay problema.
El problema surge cuando las variables de salida de las $\phi$-funciones coinciden
con las variables de entrada de otras $\phi$-funciones del mismo bloque, es decir,
cuando existe un \emph{ciclo} en el grafo de dependencias de la copia paralela.

Por ejemplo, si tras una fase de coalescencia $a_3$ se renombra a $a_1$ y $b_3$ a $b_1$,
el bloque $B_j$ quedaría:
\[
  a_1 \leftarrow \phi(b_1 : p_1,\; b_2 : p_2),\qquad
  b_1 \leftarrow \phi(a_1 : p_1,\; a_2 : p_2),
\]
y la sustitución secuencial en $p_1$:
\[
  a_1 \leftarrow b_1;\quad b_1 \leftarrow a_1
\]
es incorrecta: cuando se ejecuta $b_1 \leftarrow a_1$, la variable $a_1$ ya ha sido
sobreescrita con el valor original de $b_1$, perdiendo el valor original de $a_1$.
\end{example}

\subsection{Variables temporales frescas}

\begin{definition}[Variables temporales frescas]
\label{def:frescas}
Un conjunto de variables $\mathcal{T} = \{t_0, t_1, t_2, \ldots\} \subset \mathcal{V}$ es
\emph{fresco} con respecto a un CFG $G$ si ninguna variable de $\mathcal{T}$ aparece como
variable de salida ni como operando en ninguna instrucción o $\phi$-función de $G$.
\end{definition}

\subsection{Descripción formal del algoritmo}

\begin{definition}[Destrucción naive de SSA]
\label{def:destruccion}
Sea $G = (\mathcal{B}, \mathrm{entry})$ un CFG en forma SSA.
Sea $\mathcal{T} = \{t_0, t_1, \ldots\}$ un conjunto de variables frescas respecto a $G$,
con un contador global $k$ inicializado a $0$.

El algoritmo de \emph{destrucción naive} construye el CFG $G' = (\mathcal{B}', \mathrm{entry})$
como sigue:

\begin{enumerate}
  \item Para cada bloque $B_j \in \mathcal{B}$ con lista de $\phi$-funciones
    $\Phi_j = [\phi_1, \ldots, \phi_m]$ (donde $\phi_i.\mathrm{out} = x_i$):

    \begin{enumerate}
      \item Para cada predecesor $p \in \mathrm{pred}(B_j)$:
        \begin{enumerate}
          \item Sea $y_i = \phi_i.\mathrm{arg}(p)$ para $i = 1, \ldots, m$.

          \item \textbf{Fase 1} --- leer fuentes en temporales: añadir al final de $B_p$
            las instrucciones de copia:
            \[
              t_k \leftarrow y_1, \quad t_{k+1} \leftarrow y_2, \quad \ldots,
              \quad t_{k+m-1} \leftarrow y_m.
            \]

          \item \textbf{Fase 2} --- copiar temporales a destinos: añadir al final de $B_p$,
            tras la fase 1:
            \[
              x_1 \leftarrow t_k, \quad x_2 \leftarrow t_{k+1}, \quad \ldots,
              \quad x_m \leftarrow t_{k+m-1}.
            \]

          \item Incrementar $k \leftarrow k + m$.
        \end{enumerate}
    \end{enumerate}

  \item Para cada bloque $B_j \in \mathcal{B}$, vaciar su lista de $\phi$-funciones:
    $\Phi_j' = [\,]$.

  \item El resto de los campos de cada bloque permanecen inalterados:
    $S_j' = S_j$, $\mathrm{pred}_j' = \mathrm{pred}_j$, $\mathrm{succ}_j' = \mathrm{succ}_j$.
\end{enumerate}
\end{definition}

\begin{remark}
Las instrucciones añadidas en las fases 1 y 2 se insertan \emph{al final} del bloque predecesor
$B_p$, antes del salto al sucesor. En particular, el sucesor $\mathrm{succ}(B_p)$ no cambia.
\end{remark}

\begin{remark}
La razón por la que el uso de variables temporales frescas es correcto puede verse
intuitivamente así: en la fase 1, los valores $y_1, \ldots, y_m$ se leen en la valuación
anterior a cualquier modificación de las variables $x_1, \ldots, x_m$ (puesto que los $t_i$
son frescos, la lectura de $y_i$ no puede verse afectada por las escrituras previas de la fase
1). En la fase 2, los valores leídos se escriben en los destinos finales. El conjunto de
escrituras de la fase 2 no interfiere con las lecturas de la fase 1 porque estas ya han
terminado. Esto emula exactamente la semántica paralela de la Definición~\ref{def:ejec-phi}.
\end{remark}

\begin{example}[Resolución del ejemplo~\ref{ex:swap}]
Para el caso de intercambio $x_1 \leftarrow \phi(y_2:p)$ y $x_2 \leftarrow \phi(y_1:p)$, el
algoritmo añade en $B_p$:
\[
  \underbrace{t_0 \leftarrow y_2;\quad t_1 \leftarrow y_1}_{\text{fase 1}};\quad
  \underbrace{x_1 \leftarrow t_0;\quad x_2 \leftarrow t_1}_{\text{fase 2}}.
\]
La ejecución secuencial de estas cuatro copias produce el mismo resultado que la copia paralela
original: $x_1$ obtiene el valor (antiguo) de $y_2$ y $x_2$ obtiene el valor (antiguo) de
$y_1$.
\end{example}

% =========================================================
\section{Enunciado de corrección (sin demostración)}
% =========================================================

\subsection{Relación entre variables SSA y variables originales}

\begin{definition}[Función de des-renombramiento]
\label{def:rho}
Sea $G$ el CFG en forma SSA obtenido a partir de un programa CVM $P$. La \emph{función de
des-renombramiento} $\rho : \mathcal{V}_{\mathrm{SSA}} \to \mathcal{V}_{\mathrm{orig}}$ asigna
a cada nombre SSA $v_i$ el nombre original $x$ del que proviene.

Dado un camino de ejecución $j_0, j_1, \ldots, j_n$ en $G$ y una valuación $\sigma_n$ al
final del bloque $B_{j_n}$, definimos la \emph{valuación proyectada} sobre variables originales:
\[
  \bar\sigma_n(x) = \sigma_n(v)
  \qquad \text{donde } v \text{ es la (única) variable SSA con } \rho(v) = x
  \text{ activa al final de } B_{j_n}.
\]
\end{definition}

\begin{remark}
La noción de ``variable SSA activa al final de $B_{j_n}$'' requiere precisión formal adicional,
que se completará usando la información de vivacidad (\emph{liveness}) calculada durante la
construcción de SSA. Intuitivamente, es la última versión SSA de $x$ que ha sido definida a lo
largo del camino $j_0, \ldots, j_n$.
\end{remark}

\subsection{Equivalencia observacional}

\begin{definition}[Equivalencia de trazas]
\label{def:equiv-trazas}
Sean $G$ un CFG en forma SSA y $G' = \mathrm{Destruct}(G)$ su destrucción
(Definición~\ref{def:destruccion}). Sean
\[
  \pi  = (\sigma_0,  j_0,  \bot) \to (\sigma_1,  j_1,  j_0) \to \cdots
\]
\[
  \pi' = (\sigma_0', j_0', \bot) \to (\sigma_1', j_1', j_0') \to \cdots
\]
dos trazas de ejecución de $G$ y $G'$, respectivamente, con la misma valuación inicial de
variables originales. Decimos que $\pi$ y $\pi'$ son \emph{equivalentes} si:
\begin{enumerate}
  \item Producen el mismo camino de ejecución: $j_n = j_n'$ para todo $n$,
  \item En cada paso $n$, las valuaciones proyectadas coinciden:
    $\bar\sigma_n = \bar\sigma_n'$.
\end{enumerate}
\end{definition}

\subsection{Teorema de corrección (enunciado)}

\begin{theorem}[Corrección de la destrucción naive]
\label{thm:correctness}
Sea $P$ un programa CVM, $G$ el CFG en forma SSA obtenido por construcción estándar a partir
de $P$, y $G' = \mathrm{Destruct}(G)$ el CFG resultante de aplicar la destrucción naive
(Definición~\ref{def:destruccion}).

Para toda valuación inicial $\sigma_0$ de las variables originales de $P$, las trazas de
ejecución de $G$ y $G'$ iniciadas en $\sigma_0$ son equivalentes en el sentido de la
Definición~\ref{def:equiv-trazas}.
\end{theorem}

\subsection{Demostración del Teorema de Corrección}

% INTUICION GENERAL:
% Tenemos dos programas: G (el CFG con phis) y G' (el CFG sin phis, con copias en su lugar).
% Queremos demostrar que "hacen lo mismo". La estrategia es ir paso a paso por la ejecucion
% (bloque a bloque) y en cada paso comprobar que ambos programas estan "en el mismo sitio"
% y con "los mismos valores" en las variables que importan.
%
% La demostracion se apoya en dos lemas:
%   - Lema 1: las dos fases de copias temporales producen el mismo efecto que las phis.
%   - Lema 2: si un bloque tiene copias para varios sucesores (caso if), no se molestan entre si.
% Y luego el teorema se demuestra por induccion: si en el paso n todo va bien, en el paso n+1
% tambien.

La demostración se estructura en dos partes: un lema clave que establece que las dos fases
de copias temporales reproducen fielmente la semántica paralela de las $\phi$-funciones, y
la demostración del teorema principal por inducción sobre la longitud de la traza.

% ----------------------------------------------------------
\subsubsection{Notación para la demostración}

% INTUICION:
% Solo estamos poniendo nombres a las cosas para poder hablar de ellas:
%   - S_j son las instrucciones "de verdad" del bloque j (las mismas en G y en G').
%   - C_j^p son las copias que metimos al final del bloque p para reemplazar las phis del bloque j.
%   - S_p' es todo lo que tiene el bloque p en G': sus instrucciones originales + las copias pegadas al final.

Dado el CFG original $G = (\mathcal{B}, \mathrm{entry})$ y su destrucción
$G' = (\mathcal{B}', \mathrm{entry})$ (Definición~\ref{def:destruccion}), denotamos:

\begin{itemize}
  \item $S_j$ la secuencia de instrucciones \emph{originales} del bloque $B_j$ (igual en $G$
    y en $G'$).
  \item $C_j^p$ la secuencia de copias (fase~1 + fase~2) insertada al final de $B_p$ para las
    $\phi$-funciones de $B_j$, según la Definición~\ref{def:destruccion}.
  \item $S_p'$ la secuencia completa de instrucciones de $B_p$ en $G'$:
    \[
      S_p' = S_p \;\|\; C_{j_1}^p \;\|\; C_{j_2}^p \;\|\; \cdots
    \]
    donde $j_1, j_2, \ldots$ son los sucesores de $B_p$ que poseen $\phi$-funciones, y
    $\|$ denota concatenación.
\end{itemize}

% ----------------------------------------------------------
\subsubsection{Lema clave: corrección de las dos fases de copias}

% INTUICION DEL LEMA 1:
% La phi dice "lee todo a la vez y escribe todo a la vez" (paralelo).
% Nosotros no podemos hacer eso con instrucciones normales (que son secuenciales).
% Lo que hacemos es:
%   Fase 1: guardar cada valor original en un temporal (tmp_0 = y1, tmp_1 = y2, ...)
%   Fase 2: copiar cada temporal al destino final (x1 = tmp_0, x2 = tmp_1, ...)
%
% Esto funciona porque:
%   - En la Fase 1, escribir en un temporal NUNCA machaca ningun y_i,
%     porque los temporales son nombres nuevos que no existen en el programa original.
%   - En la Fase 2, escribir en un destino x_j NUNCA machaca ningun temporal,
%     por la misma razon: son nombres completamente distintos.
%
% Asi que al final, cada x_i acaba con el valor original de y_i, exactamente como la phi.

\begin{lemma}[Las copias temporales emulan la ejecución paralela]
\label{lem:two-phase}
Sea $B_j$ un bloque con $\Phi_j = [\phi_1, \ldots, \phi_m]$, donde
$\phi_i.\mathrm{out} = x_i$ y $\phi_i.\mathrm{arg}(p) = y_i$ para un predecesor $p$.
Sean $t_k, t_{k+1}, \ldots, t_{k+m-1}$ variables frescas (Definición~\ref{def:frescas}).

Sea $\sigma$ una valuación cualquiera. Entonces la ejecución secuencial de las copias de
fase~1 y fase~2:
\[
  C_j^p = [\,t_k \leftarrow y_1,\; \ldots,\; t_{k+m-1} \leftarrow y_m,\;
            x_1 \leftarrow t_k,\; \ldots,\; x_m \leftarrow t_{k+m-1}\,]
\]
satisface:
\begin{enumerate}
  \item $\llbracket C_j^p \rrbracket_\sigma(x_i)
        = \llbracket \Phi_j \rrbracket_{\sigma,\,p}(x_i)
        = \sigma(y_i)$
    \quad para todo $i \in \{1, \ldots, m\}$.

  \item $\llbracket C_j^p \rrbracket_\sigma(v) = \sigma(v)$
    \quad para todo $v \notin \{x_1, \ldots, x_m,\; t_k, \ldots, t_{k+m-1}\}$.
\end{enumerate}
\end{lemma}

\begin{proof}

% INTUICION DE LA PRUEBA:
% Vamos a seguir la ejecucion instruccion por instruccion.
% sigma_0 es el estado ANTES de ejecutar nada.
% sigma_1 es el estado despues de ejecutar "tmp_0 = y1".
% sigma_2 es el estado despues de ejecutar tambien "tmp_1 = y2".
% ... y asi hasta sigma_m (fin de Fase 1).
% Luego seguimos con la Fase 2 hasta llegar al estado final sigma*.

Definimos $\sigma_0 = \sigma$ y, para $i = 1, \ldots, m$:
\[
  \sigma_i = \sigma_{i-1}[t_{k+i-1} \mapsto \sigma_{i-1}(y_i)].
\]
Estas son las valuaciones intermedias tras cada instrucción de la fase~1.

\medskip
% INTUICION DE LA FASE 1:
% Queremos demostrar que cada temporal acaba con el valor original de su fuente.
% Es decir: tmp_i guarda el valor que tenia y_i ANTES de que empezaramos a tocar nada.
% Esto es asi porque:
%   - Cuando escribimos en tmp_i, NO estamos modificando ningun y_j (los temporales
%     son nombres completamente nuevos, nunca coinciden con variables del programa).
%   - Tampoco estamos modificando ningun temporal anterior (cada tmp tiene nombre distinto).
% Entonces la lectura de y_i siempre da el valor original.
\textbf{Propiedad de la fase~1.} Para todo $i \in \{1, \ldots, m\}$:
\[
  \sigma_i(t_{k+i-1}) = \sigma(y_i).
\]

Esto se demuestra por inducción sobre $i$. Para $i = 1$:
$\sigma_1(t_k) = \sigma_0(y_1) = \sigma(y_1)$.
Para el paso inductivo, basta observar que la escritura de $t_{k+i-1}$ no modifica
ninguna variable $y_l$ (con $l \leq m$) ni ningún temporal previo $t_{k+l-1}$ (con $l < i$),
porque:
\begin{itemize}
  \item Los temporales $t_{k+i-1}$ son \emph{frescos}: $t_{k+i-1} \neq y_l$ para todo $l$
    (ningún temporal aparece como operando en $G$, y los $y_l$ son operandos de $\phi$-funciones
    de~$G$).
  \item Todos los temporales son distintos entre sí por construcción del contador $k$.
\end{itemize}

Por lo tanto, al final de la fase~1 (tras $m$ instrucciones), la valuación $\sigma_m$ satisface:
\begin{align}
  \sigma_m(t_{k+i-1}) &= \sigma(y_i) \quad \text{para } i = 1, \ldots, m, \label{eq:phase1-temps}\\
  \sigma_m(v) &= \sigma(v) \quad \text{para } v \notin \{t_k, \ldots, t_{k+m-1}\}. \label{eq:phase1-rest}
\end{align}

\medskip
% INTUICION DE LA FASE 2:
% Ahora los temporales ya tienen guardados los valores originales.
% Vamos asignando: x1 = tmp_0, x2 = tmp_1, ...
% La pregunta clave: cuando hacemos x2 = tmp_1, no se habra machacado tmp_1?
% No, porque escribimos en x_j (que es una variable SSA del programa), y los temporales
% son nombres frescos. x_j nunca es igual a tmp_i. Asi que leer tmp_i siempre da bien.
\textbf{Propiedad de la fase~2.} Definimos $\sigma_m^{(0)} = \sigma_m$ y, para $i = 1, \ldots, m$:
\[
  \sigma_m^{(i)} = \sigma_m^{(i-1)}[x_i \mapsto \sigma_m^{(i-1)}(t_{k+i-1})].
\]
La valuación final es $\sigma^* = \sigma_m^{(m)} = \llbracket C_j^p \rrbracket_\sigma$.

Para todo $i$, se tiene que $\sigma_m^{(i-1)}(t_{k+i-1}) = \sigma_m(t_{k+i-1}) = \sigma(y_i)$.
Esto se debe a que la escritura $x_l \leftarrow t_{k+l-1}$ (para $l < i$) modifica $x_l$, pero
$x_l \neq t_{k+i-1}$ ya que los temporales son frescos (no son variables SSA de~$G$, mientras
que los~$x_l$ sí lo son). Por tanto:
\[
  \sigma^*(x_i) = \sigma(y_i) \quad \text{para todo } i \in \{1, \ldots, m\}.
\]

Además, para toda variable $v \notin \{x_1, \ldots, x_m,\, t_k, \ldots, t_{k+m-1}\}$,
ninguna de las $2m$ instrucciones de $C_j^p$ modifica $v$, luego $\sigma^*(v) = \sigma(v)$.

\medskip
Comparando con la Definición~\ref{def:ejec-phi}:
$\llbracket \Phi_j \rrbracket_{\sigma,\,p}(x_i) = \sigma(y_i^p) = \sigma(y_i)$,
concluimos que $\sigma^*(x_i) = \llbracket \Phi_j \rrbracket_{\sigma,\,p}(x_i)$ para todo~$i$.
\end{proof}

% ----------------------------------------------------------
\subsubsection{Propiedad estructural: a lo sumo un sucesor con $\phi$-funciones}

% INTUICION:
% Antes teniamos un Lema de "no interferencia" para el caso en que un bloque
% tuviera dos sucesores condicionales ambos con phis. Pero al revisar la
% construccion de SSA, resulta que eso NUNCA pasa:
%   - Los saltos condicionales solo se crean en handle_if.
%   - Alli, los dos sucesores inmediatos (then_b y else_b) son bloques
%     recien creados, cada uno con UN SOLO predecesor (el bloque condicional).
%   - Con un solo predecesor, read_recursive nunca inserta phis.
%   - El unico caso en que un sucesor del bloque condicional puede tener phis
%     es el if SIN else, donde else_b = join. El join tiene 2+ predecesores,
%     pero then_b sigue teniendo solo 1.
%   - Por tanto, como mucho UN sucesor tiene phis. Nunca los dos.
%
% Esto simplifica la demostracion: no hace falta demostrar que las copias
% de distintos sucesores no interfieren, porque esa situacion no existe.

\begin{proposition}[A lo sumo un sucesor con $\phi$-funciones]
\label{prop:at-most-one-phi-succ}
Sea $G = (\mathcal{B}, \mathrm{entry})$ un CFG en forma SSA construido por el algoritmo
de construcción estándar (Braun et al.). Para todo bloque $B_p$ con sucesor condicional
$\mathrm{succ}(B_p) = \mathrm{cond}(a, j_1, j_2)$, a lo sumo uno de los bloques $B_{j_1}$,
$B_{j_2}$ tiene una lista de $\phi$-funciones no vacía.
\end{proposition}

\begin{proof}
En la construcción de SSA implementada, los saltos condicionales se generan únicamente en
el manejo de sentencias \texttt{if-then-else}. Analizamos los dos casos posibles:

\medskip
\emph{Caso 1: \texttt{if} con \texttt{else}.}
Se crean tres bloques nuevos: $B_{\mathrm{then}}$, $B_{\mathrm{else}}$ y $B_{\mathrm{join}}$.
El bloque actual $B_p$ recibe el sucesor condicional hacia $B_{\mathrm{then}}$ y $B_{\mathrm{else}}$.
Ambos son bloques recién creados, y su único predecesor es~$B_p$:
$|\mathrm{pred}(B_{\mathrm{then}})| = |\mathrm{pred}(B_{\mathrm{else}})| = 1$.

Las $\phi$-funciones solo se insertan cuando un bloque tiene dos o más predecesores
(en el algoritmo, \texttt{read\_recursive} solo crea una $\phi$ cuando
$|\mathrm{pred}(B)| \geq 2$). Por tanto, ni $B_{\mathrm{then}}$ ni $B_{\mathrm{else}}$
contienen $\phi$-funciones. Los sucesores directos de $B_p$ tienen ambos
$\Phi = [\,]$.

\medskip
\emph{Caso 2: \texttt{if} sin \texttt{else}.}
Se crean dos bloques nuevos: $B_{\mathrm{then}}$ y $B_{\mathrm{join}}$.
El bloque actual $B_p$ recibe el sucesor condicional hacia $B_{\mathrm{then}}$ (rama verdadera) y
$B_{\mathrm{join}}$ (rama falsa).

$B_{\mathrm{then}}$ tiene un único predecesor ($B_p$), luego $\Phi_{\mathrm{then}} = [\,]$.
$B_{\mathrm{join}}$ recibe como predecesores a $B_p$ (rama falsa) y al bloque final del cuerpo
del \texttt{then}, luego $|\mathrm{pred}(B_{\mathrm{join}})| \geq 2$ y puede contener
$\phi$-funciones.

En este caso, exactamente uno de los dos sucesores ($B_{\mathrm{join}}$) puede tener
$\phi$-funciones, y el otro ($B_{\mathrm{then}}$) no.

\medskip
En ambos casos, a lo sumo un sucesor directo de $B_p$ tiene $\Phi \neq [\,]$.
\end{proof}

\begin{remark}
Esta propiedad es específica de la construcción de SSA implementada. En un CFG SSA
genérico (por ejemplo, tras transformaciones arbitrarias), podría darse el caso de un
bloque con dos sucesores condicionales ambos con $\phi$-funciones. Para completitud,
enunciamos a continuación el lema de no interferencia que cubriría ese caso general,
aunque no es necesario para la corrección de nuestra implementación concreta.
\end{remark}

\begin{lemma}[No interferencia --- caso general]
\label{lem:no-interference}
Sea $B_p$ un bloque con sucesores $B_{j_1}$ y $B_{j_2}$ (caso condicional), ambos con
$\phi$-funciones. Las copias $C_{j_1}^p$ y $C_{j_2}^p$ insertadas en $B_p$ no interfieren
entre sí: la ejecución secuencial
$\llbracket S_p \;\|\; C_{j_1}^p \;\|\; C_{j_2}^p \rrbracket_\sigma$
produce los mismos valores para las variables de salida de $\Phi_{j_1}$ y $\Phi_{j_2}$ que
si se ejecutase cada grupo de copias por separado sobre $\llbracket S_p \rrbracket_\sigma$.
\end{lemma}

\begin{proof}
Sean $\{x_1^{(1)}, \ldots, x_{m_1}^{(1)}\}$ las variables de salida de $\Phi_{j_1}$ y
$\{x_1^{(2)}, \ldots, x_{m_2}^{(2)}\}$ las de $\Phi_{j_2}$. Sean
$\{t_k, \ldots, t_{k+m_1-1}\}$ los temporales de $C_{j_1}^p$ y
$\{t_{k+m_1}, \ldots, t_{k+m_1+m_2-1}\}$ los de $C_{j_2}^p$.

Los siguientes conjuntos son \emph{mutuamente disjuntos}:
\begin{itemize}
  \item Las variables de salida de $\Phi_{j_1}$ y las de $\Phi_{j_2}$ son disjuntas por la
    propiedad SSA (Definición~\ref{def:ssa}): cada variable se define a lo sumo una vez.
  \item Todos los temporales son distintos entre sí (contador global $k$) y distintos de toda
    variable SSA (frescura, Definición~\ref{def:frescas}).
\end{itemize}

Por tanto, las escrituras de $C_{j_1}^p$ (a temporales $t_k, \ldots, t_{k+m_1-1}$ y a variables
$x^{(1)}_i$) no modifican ninguna variable leída ni escrita por $C_{j_2}^p$ (que lee
$y^{(2)}_i$ y escribe temporales $t_{k+m_1}, \ldots$ y variables $x^{(2)}_i$), y viceversa.

En particular, las fuentes $y_i^{(l)}$ de cada grupo de copias son variables definidas en
bloques anteriores de~$G$, y por tanto no son ni temporales ni variables de salida de phis
de otros bloques. Así, las escrituras de un grupo no afectan las lecturas del otro.
\end{proof}

% ----------------------------------------------------------
\subsubsection{Demostración del Teorema~\ref{thm:correctness}}

% INTUICION DEL TEOREMA PRINCIPAL:
% Imaginemos que G y G' son dos "maquinas" que ejecutan el programa bloque a bloque.
% Queremos demostrar que siempre van al mismo bloque y producen los mismos valores.
%
% Lo hacemos por induccion: si en el paso n ambas maquinas estan en el mismo bloque
% con los mismos valores relevantes, entonces en el paso n+1 tambien.
%
% El "invariante" I(n) dice exactamente eso:
%   (I1) Estan en el mismo bloque (mismo camino).
%   (I2) Las variables que importan para ejecutar ese bloque tienen el mismo valor.
%        En G, despues de ejecutar las phis del bloque.
%        En G', directamente al entrar (porque las copias ya se hicieron en el bloque anterior).

\begin{proof}[Demostración del Teorema~\ref{thm:correctness}]
La demostración procede por inducción fuerte sobre el índice $n$ de la traza.

\medskip
\textbf{Invariante $I(n)$.} Sea
$\pi = (\sigma_0, j_0, \bot) \to (\sigma_1, j_1, p_1) \to \cdots$
la traza de $G$ y
$\pi' = (\sigma_0', j_0', \bot) \to (\sigma_1', j_1', p_1') \to \cdots$
la traza de $G'$, ambas con la misma valuación inicial de variables originales.

Decimos que $I(n)$ se cumple si:
\begin{enumerate}
  \item[\textbf{(I1)}] $j_n = j_n'$ y $p_n = p_n'$ \quad (mismo camino),
  \item[\textbf{(I2)}] para toda variable $v$ que sea operando de alguna instrucción
    de~$S_{j_n}$ o de la condición de $\mathrm{succ}(B_{j_n})$, se tiene que la valuación al
    inicio de la ejecución de $S_{j_n}$ es la misma en $G$ y en $G'$. Formalmente:
    \[
      \llbracket \Phi_{j_n} \rrbracket_{\sigma_n^-,\, p_n}(v)
      \;=\; \sigma_n'^{-}(v),
    \]
    donde $\sigma_n^-$ y $\sigma_n'^{-}$ son las valuaciones al \emph{entrar} en el bloque
    $j_n$ en $G$ y $G'$ respectivamente (es decir, antes de ejecutar $\Phi_{j_n}$ y $S_{j_n}$
    en $G$, y antes de ejecutar $S_{j_n}$ en $G'$, que no tiene phis).
\end{enumerate}

\medskip
% CASO BASE: Ambos programas empiezan en el mismo sitio (bloque entry) con los mismos valores.
% El bloque de entrada no tiene predecesores, asi que no tiene phis. No hay diferencia
% entre G y G' aqui. Trivial.
\textbf{Caso base ($n = 0$).}
En ambos CFGs, $j_0 = j_0' = \mathrm{entry}$ y $p_0 = p_0' = \bot$.
El bloque de entrada no tiene predecesores, luego $\Phi_{\mathrm{entry}} = [\,]$
(no hay $\phi$-funciones sin predecesores). Así
$\llbracket \Phi_{\mathrm{entry}} \rrbracket_{\sigma_0, \bot} = \sigma_0$.
En $G'$, $\sigma_0'^{-} = \sigma_0' = \sigma_0$ (misma valuación inicial).
El invariante $I(0)$ se satisface trivialmente.

\medskip
% PASO INDUCTIVO:
% Suponemos que I(n) vale (estamos en el mismo bloque, mismos valores).
% Ahora ejecutamos el bloque actual y vemos que pasa.
%
% Paso 1: Las instrucciones originales S_{j_n} son las MISMAS en G y G'.
%         Y por I(n) reciben los mismos valores de entrada.
%         Por tanto producen los mismos resultados. Llamamos sigma_mid a ese estado comun.
\textbf{Paso inductivo.} Supongamos que $I(n)$ se cumple. Probamos $I(n+1)$.

Por la hipótesis \textbf{(I2)}, las instrucciones de $S_{j_n}$ reciben los mismos valores de
entrada en $G$ y en $G'$. Puesto que $S_{j_n}$ es idéntico en ambos CFGs
(la destrucción no modifica las instrucciones originales), la ejecución secuencial
(Definición~\ref{def:ejec-secuencia}) produce la misma valuación para todas las variables
definidas en $S_{j_n}$ y preserva el valor de las demás. Llamemos $\sigma_{\mathrm{mid}}$
a esta valuación común tras ejecutar $S_{j_n}$ (en $G$) y $S_{j_n}$ (la parte original, en $G'$).

\medskip
% Paso 2: Como los valores son iguales y el sucesor (la condicion del if o el goto) no ha
%         cambiado, ambos programas van al MISMO bloque siguiente.
%         Esto nos da (I1) para n+1.
\textbf{Determinación del siguiente bloque.}
En $G$: $j_{n+1} = \mathrm{next}(j_n, \sigma_{\mathrm{mid}})$.
En $G'$: el sucesor $\mathrm{succ}(B_{j_n})$ es el mismo (no se modifica), y la condición
se evalúa sobre $\sigma_{\mathrm{mid}}$ (la condición solo depende de variables de $S_{j_n}$
o anteriores, que son iguales). Luego $j_{n+1}' = j_{n+1}$ y $p_{n+1}' = j_n = j_n' = p_{n+1}$.
Se cumple \textbf{(I1)} para $n+1$.

\medskip
% Paso 3: Ahora viene la parte clave. En G', despues de las instrucciones originales,
%         se ejecutan las copias que metimos para reemplazar las phis del bloque siguiente.
%
%         En G, las phis del bloque siguiente se ejecutarian al ENTRAR en ese bloque.
%         En G', las copias se ejecutan al SALIR del bloque actual.
%         Pero el efecto es el mismo (Lema 1).
%
%         Si ademas hay copias "extra" para phis de OTRO sucesor (caso if con dos ramas),
%         no pasa nada: no interfieren (Lema 2).
\textbf{Efecto de las copias en $G'$.}
En $G'$, tras ejecutar $S_{j_n}$ (dando $\sigma_{\mathrm{mid}}$), se ejecutan las copias
$C_{j_{n+1}}^{j_n}$ (y posiblemente $C_{j'}^{j_n}$ para otros sucesores $j'$ de $B_{j_n}$
que tengan phis). Sea $\sigma_{\mathrm{fin}}'$ la valuación tras ejecutar todas las copias.

Por el Lema~\ref{lem:two-phase}, aplicado con $\sigma = \sigma_{\mathrm{mid}}$:
\[
  \sigma_{\mathrm{fin}}'(x_i) = \sigma_{\mathrm{mid}}(y_i)
  \quad \text{para cada } \phi_i \in \Phi_{j_{n+1}} \text{ con } \phi_i.\mathrm{arg}(j_n) = y_i.
\]

Por el Lema~\ref{lem:no-interference}, las copias para otros sucesores $j' \neq j_{n+1}$
no modifican las variables $y_i$ (fuentes de $\Phi_{j_{n+1}}$), ni las variables de salida
$x_i$ de $\Phi_{j_{n+1}}$, ni los temporales usados por $C_{j_{n+1}}^{j_n}$.

Para toda variable $v$ que no es variable de salida de ningún $\Phi$ ni temporal:
$\sigma_{\mathrm{fin}}'(v) = \sigma_{\mathrm{mid}}(v)$.

\medskip
% Paso 4: Ahora verificamos (I2) para n+1.
%         Cogemos cualquier variable v que se use en las instrucciones del bloque siguiente.
%         Hay dos casos:
%
%         Caso 1: v NO es una variable de salida de ninguna phi del bloque siguiente.
%                 Entonces ni las phis (en G) ni las copias (en G') la tocan.
%                 Su valor viene de antes, y es el mismo: sigma_mid(v).
%
%         Caso 2: v SI es una variable de salida de una phi (v = x_i).
%                 En G: la phi le asigna el valor sigma_mid(y_i).
%                 En G': la copia (Lema 1) le asigna el valor sigma_mid(y_i).
%                 Mismo resultado.
\textbf{Verificación de \textbf{(I2)} para $n+1$.}
Sea $v$ un operando de $S_{j_{n+1}}$ o de la condición de $\mathrm{succ}(B_{j_{n+1}})$.

\emph{Caso 1:} $v$ no es variable de salida de ninguna $\phi$-función de $\Phi_{j_{n+1}}$.
\begin{itemize}
  \item En $G$: $\llbracket \Phi_{j_{n+1}} \rrbracket_{\sigma_{\mathrm{mid}}, j_n}(v)
    = \sigma_{\mathrm{mid}}(v)$ \quad (la ejecución paralela no modifica $v$).
  \item En $G'$: $\sigma_{\mathrm{fin}}'(v) = \sigma_{\mathrm{mid}}(v)$ \quad
    (las copias no escriben en $v$: los destinos son variables de salida de phis y temporales,
    ninguno de los cuales es $v$).
  \item Ambos coinciden. \checkmark
\end{itemize}

\emph{Caso 2:} $v = x_i$ para alguna $\phi_i \in \Phi_{j_{n+1}}$ con
$\phi_i.\mathrm{arg}(j_n) = y_i$.
\begin{itemize}
  \item En $G$: $\llbracket \Phi_{j_{n+1}} \rrbracket_{\sigma_{\mathrm{mid}}, j_n}(x_i)
    = \sigma_{\mathrm{mid}}(y_i)$ \quad (Definición~\ref{def:ejec-phi}).
  \item En $G'$: $\sigma_{\mathrm{fin}}'(x_i) = \sigma_{\mathrm{mid}}(y_i)$ \quad
    (Lema~\ref{lem:two-phase}).
  \item Ambos coinciden. \checkmark
\end{itemize}

\medskip
En ambos casos \textbf{(I2)} se cumple para $n+1$.

\medskip
% CONCLUSION: Como (I1) e (I2) se cumplen para n+1, podemos seguir indefinidamente.
% En cada paso ambos programas estan en el mismo sitio con los mismos valores.
% Por tanto las trazas son equivalentes.
\textbf{Equivalencia de valuaciones proyectadas.}
Dado que $I(n)$ implica que $S_{j_n}$ se ejecuta con los mismos operandos en $G$ y $G'$,
las variables definidas por $S_{j_n}$ reciben los mismos valores. Las variables ``activas''
al final de $B_{j_n}$ (en el sentido de la Definición~\ref{def:rho}) son aquellas definidas
en $S_{j_n}$ o heredadas de bloques anteriores; en ambos casos sus valores coinciden.
Por tanto, $\bar\sigma_n = \bar\sigma_n'$ (las valuaciones proyectadas son iguales).

\medskip
Por inducción, $I(n)$ se cumple para todo $n \geq 0$, lo que establece la equivalencia
de trazas en el sentido de la Definición~\ref{def:equiv-trazas}.
\end{proof}

\end{document}
